% The five stages of the pipeline, with romanization as a side branch.
\begin{tikzpicture}[
  stage/.style={draw=blue!55!black, rounded corners=4pt, fill=blue!7, line width=0.8pt,
                minimum width=10.9cm, minimum height=2.0cm, inner sep=6pt,
                execute at begin node={\hyphenpenalty 10000\relax}},
  rule/.style={stage, draw=orange!70!black, fill=orange!10},
  io/.style={draw=black!60, rounded corners=10pt, fill=black!5, line width=0.8pt,
             text width=5.2cm, minimum height=1.0cm, align=center, font=\small},
  data/.style={font=\footnotesize\itshape, text=black!70, anchor=west, align=left, text width=6.6cm},
  arr/.style={-{Stealth[length=3mm]}, line width=1pt},
  title/.style={font=\bfseries\large}
]
% ---- nodes ----
\node[io] (in) at (0,0) {\faIcon{microphone}\ \textbf{Input:} Bengali speech clip\\ \footnotesize (the original recording)};

\node[stage] (asr) at (0,-2.4) {\stagecontent{blue!55!black}{\faIcon{file-alt}}{1. Transcription}
  {Converts the Bengali speech into Bengali text written in native script.}};

% Romanization forms a side branch; the transcript itself goes straight to Stage 3.
\node[rule, minimum width=6.5cm] (rom) at (2.2,-5.2) {\stagecontent[4.6cm]{orange!70!black}{\faIcon{font}}{2. Romanization}
  {Rewrites each word in Latin letters for the aligner.}};

\node[stage] (pro) at (0,-8.0) {\stagecontent{blue!55!black}{\faIcon{wave-square}}{3. Alignment and prosody marking}
  {Finds when each word starts and ends, then inserts \texttt{,} for short pauses and \texttt{|} for
  long pauses; also measures loudness and elongated words.}};

\node[stage] (tr) at (0,-10.8) {\stagecontent{blue!55!black}{\faIcon{language}}{4. Translation}
  {Translates the pause-marked Bengali text into English.}};

\node[stage] (tts) at (0,-13.6) {\stagecontent{blue!55!black}{\faIcon{volume-up}}{5. Voice-cloned speech synthesis}
  {Speaks the English text in the original speaker's voice, copied zero-shot from the input clip.}};

\node[io] (out) at (0,-16.1) {\faIcon{headphones}\ \textbf{Output:} English speech\\ \footnotesize in the original speaker's voice};

% ---- main flow with data labels ----
\draw[arr] (in)  -- node[data, xshift=2mm]{audio} (asr);
\draw[arr] (asr.south -| rom) -- node[data, xshift=2mm, text width=3cm]{Bengali words} (rom);
\draw[arr] (rom) -- node[data, xshift=2mm, text width=3cm]{romanized words} (rom.south |- pro.north);
\draw[arr] (asr.south -| -3.8,0) -- node[data, xshift=2mm, text width=2.4cm]{Bengali text, no punctuation} (pro.north -| -3.8,0);
\draw[arr] (pro) -- node[data, xshift=2mm]{Bengali sentence with pause markers} (tr);
\draw[arr] (tr)  -- node[data, xshift=2mm]{English translation} (tts);
\draw[arr] (tts) -- node[data, xshift=2mm]{synthesized audio} (out);

% the aligner also needs the audio
\draw[arr] (in.east) -- (6.3,0) |- (pro.east);
\node[font=\footnotesize\itshape, text=black!70, rotate=90, anchor=south] at (6.3,-4.0) {audio};

% ---- side inputs ----
\draw[arr, dashed, black!55] (in.west) -- (-6.6,0) |- ([yshift=4mm]tts.west);
\node[font=\footnotesize\itshape, text=black!65, rotate=90, anchor=south] at (-6.6,-6.8)
  {original clip reused as the voice reference};

\draw[arr, dashed, black!55] ([yshift=-4mm]pro.west) -- (-5.8,-8.4) |- ([yshift=-4mm]tts.west);
\node[font=\footnotesize\itshape, text=black!65, rotate=90, anchor=south] at (-5.8,-11.0)
  {marked Bengali as transcript};
\end{tikzpicture}
